\documentclass[letterpaper,12pt]{article}
\usepackage{amsmath, amssymb, geometry}
\geometry{margin=1in}
\setlength{\parindent}{0pt} % no indentation
\begin{document}

%INFO
Christian DiPietrantonio \\
ME 5130: Advanced Heat and Mass Transfer \\
Homework 01: Heat Conduction \\
\today

\vspace{6pt}

\textbf{1--2, Problem Statement:}
Verify that $\nabla T$ and $\nabla \cdot \mathbf{q}$ in the spherical coordinate system $\left(r, \phi, \theta \right)$ are given as:

\begin{equation}
    \nabla T = \mathbf{\hat{u}_r} \frac{\partial T}{\partial r} + \mathbf{\hat{u}_{\phi}} \frac{1}{r \sin \theta} \frac{\partial T}{\partial \phi} + \mathbf{\hat{u}_{\theta}} \frac{1}{r} \frac{\partial T} {\partial \theta}
\end{equation}

\begin{equation}
    \nabla \cdot \mathbf{q} = \frac{1}{r^2} \frac{\partial}{\partial r} \left(r^2 q_r \right) + \frac{1}{r \sin \theta} \frac{\partial q_{\phi}}{\partial \phi} + \frac{1}{r \sin \theta} \frac{\partial}{\partial \theta} \left( q_{\theta} \sin \theta \right)
\end{equation}

\vspace{6pt}

\textbf{1--2, Solution:}
The transformation from spherical to cartesioan coordinates is:

\begin{equation}
\begin{aligned}
    x &= r\sin{\theta}\cos{\phi} \\
    y &= r\sin{\theta}\sin{\phi} \\
    z &= r\cos{\theta}
\end{aligned}
\end{equation}

and the following can be defined:

\begin{equation}
\begin{aligned}
    u_1 &= r \\
    u_2 &= \phi \\
    u_3 &= \theta
\end{aligned}
\end{equation}

Now, from \"Oz\i \c{s}\i k, we have Eq. (1--19):

\begin{equation}
    \nabla T = \mathbf{\hat{u}_1} \frac{1}{a_1} \frac{\partial T}{\partial u_1} + \mathbf{\hat{u}_2} \frac{1}{a_2} \frac{\partial T}{\partial u_2} + \mathbf{\hat{u}_3} \frac{1}{a_3} \frac{\partial T}{\partial u_3}
\end{equation}

Eq. (1--22a):

\begin{equation}
    \nabla \cdot \mathbf{q} = \frac{1}{a} \left[ \frac{\partial}{\partial u_1} \left( \frac{a}{a_1} q_1 \right) + \frac{\partial}{\partial u_2} \left( \frac{a}{a_2} q_2 \right) + \frac{\partial}{\partial u_3} \left( \frac{a}{a_3} q_3 \right) \right]
\end{equation}

Eq. (1--22b):

\begin{equation}
    a = a_1 a_2 a_3
\end{equation}

and Eq (1--18):

\begin{equation}
    a^{2}_{i} = \left(\frac{\partial x}{\partial u_i} \right)^2 + \left(\frac{\partial y}{\partial u_i} \right)^2 + \left(\frac{\partial z}{\partial u_i} \right)^2, \qquad i = 1, 2, 3
\end{equation}

We can now compute the partial derivatives needed to find the scale factors:

\begin{equation}
    \frac{\partial x}{\partial r} = \sin{\theta} \cos{\phi}
\end{equation}

\begin{equation}
    \frac{\partial y}{\partial r} = \sin{\theta} \sin{\phi}
\end{equation}

\begin{equation}
    \frac{\partial z}{\partial r} = \cos{\theta}
\end{equation}

\begin{equation}
    \frac{\partial x}{\partial \phi} = -r \sin{\theta} \sin{\phi}
\end{equation}

\begin{equation}
    \frac{\partial y}{\partial \phi} = r \sin{\theta} \cos{\phi}
\end{equation}

\begin{equation}
    \frac{\partial z}{\partial \phi} = 0
\end{equation}

\begin{equation}
    \frac{\partial x}{\partial \theta} = r \cos{\theta} \cos{\phi}
\end{equation}

\begin{equation}
    \frac{\partial y}{\partial \theta} = r \cos{\theta} \sin{\phi}
\end{equation}

\begin{equation}
    \frac{\partial z}{\partial \theta} = -r \sin{\theta}
\end{equation}

We can now evaluate:

\begin{equation}
\begin{aligned}
     a^{2}_{r} &= \sin^2{\theta} \cos^2{\phi} + \sin^2{\theta} \sin^2{\phi} + \cos^2{\theta} \\
               &= \sin^2{\theta} (\cos^2{\phi} + \sin^2{\theta}) + \cos^2{\theta} \\
               &= \sin^2{\theta} + \cos^2{\theta} \\
               &= 1 \\
     a_1 = a_r &= 1
\end{aligned}
\end{equation}

\begin{equation}
\begin{aligned}
    a^{2}_{\phi}   &= r^2 \sin^2{\theta} \sin^2{\phi} + r^2 \sin^2{\theta} \cos^2{\phi} \\
                   &= r^2 \sin^2{\theta} (\sin^2{\phi} + \cos^2{\phi}) \\
                   &= r^2 \sin^2{\theta} \\
    a_2 = a_{\phi} &= r \sin{\theta}
\end{aligned}
\end{equation}

\begin{equation}
\begin{aligned}
    a^{2}_{\theta}   &= r^2 \cos^2{\theta} \cos^2{\phi} + r^2 \cos^2{\theta} \sin^2{\phi} + r^2 \sin^2{\theta} \\
                     &= r^2 \cos^2{\theta} (\cos^2{\phi} + \sin^2{\phi}) + r^2 \sin^2{\theta} \\
                     &= r^2 \cos^2{\theta} + r^2 \sin^2{\theta} \\
                     &= r^2 (\cos^2{\theta} + \sin^2{\theta}) \\
                     &= r^2 \\
    a_3 = a_{\theta} &= r       
\end{aligned}
\end{equation}

We should note that since spherical coordinates ($r, \phi, \theta$) form an orthogonal coordinate system, we can use equations (1--19) and (1--22). Nevertheless, we can now re-write equation (1--19):

\begin{equation}
\boxed{
    \nabla T = \mathbf{\hat{u}_r} \frac{\partial T}{\partial r} + \mathbf{\hat{u}_{\phi}} \frac{1}{r \sin{\theta}} \frac{\partial T}{\partial \phi} + \mathbf{\hat{u}_{\theta}} \frac{1}{r} \frac{\partial T}{\partial \theta}
}
\end{equation}

We can also now write:

\begin{equation}
    a = a_1 a_2 a_3 = a_r a_{phi} a_{\theta} = r^2 \sin{\theta}
\end{equation}

Therefore:

\begin{equation}
    \frac{a}{a_r} = r^2 \sin{theta}
\end{equation}

\begin{equation}
    \frac{a}{a_{\phi}} = r
\end{equation}'

\begin{equation}
    \frac{a}{a_{\theta}} = r \sin{theta}
\end{equation}

Therefore, equation (1--22a) can be re-written as:

\begin{equation}
    \nabla \cdot \mathbf{q} = \frac{1}{r^2 \sin{\theta}} \left[ \frac{\partial}{\partial r} \left( r^2 \sin{\theta} q_r \right) + \frac{\partial}{\partial \phi} \left( r q_{\phi} \right) + \frac{\partial}{\partial \theta} \left( r \sin{\theta} q_{\theta} \right) \right]
\end{equation}

Which can be simplified to:

\begin{equation}
\boxed{
    \nabla \cdot \mathbf{q} = \frac{1}{r^2} \frac{\partial}{\partial r} \left( r^2 q_r \right) + \frac{1}{r \sin{\theta}}\frac{\partial q_{\phi}}{\partial \phi} + \frac{1}{r \sin{\theta}}\frac{\partial}{\partial \theta} \left(\sin{\theta} q_{\theta} \right)
}
\end{equation}



\vspace{6pt}

\textbf{1--12, Problem Statement:}
Set up the mathematical formulation of the following head conduction problems:

\begin{enumerate}
    \item A slab in $0 \leq x \leq L$ is initially at a temperature $F\left( x \right)$. For times $t > 0$, the boundary at $x = 0$ is kept insulated and the boundary at $x = L$ dissipates heat by convection into a medium at zero temperature.
    \item A semiinfinite region $0 \leq x < \infty$ is initially at a temperature $F\left( x \right)$. For times $t > 0$, heat is generated in the medium at a constant rate of $g_0 W/m^3$, while the boundary at x = 0 is kept at zero temperature.
    \item A solid cylinder $0 \leq r \leq b$ is initially at a temperature $F\left( r \right)$. For times $t > 0$, heat is generated in the medium at a rate of $g\left( r \right)$, $W/m^3$, while the boundary at $r = b$ dissipates heat by convection into a medium at zero temperature.
    \item A solid sphere $0 \leq r \leq b$ is initially at temperature $F\left( r \right)$. For times $t > 0$, heat is generated in the medoum at a rate of $g\left( r \right)$, $W/m^3$, while the boundary at $r = b$ is kept at a uniform temperature $T_0$.
\end{enumerate}

\vspace{6pt}

\textbf{1--12, Solution:}

\vspace{6pt}

\textbf{1--15, Problem Statement:}
A long cylindrical iron bar of diameter $D = 5\, \text{cm}$, initially at temperature $T_0 = 650\, {}^{\circ} \text{C}$, is exposed to an air stream at $T_{\infty} = 50\, {}^{\circ} \text{C}$. The heat transfer coefficient between the air stream and the surface of the bar is $h = 80\, \text{W}\,\text{m}^{-2}\,{}^\circ\text{C}^{-1}$. Thermophysical properties may be taken as $\rho = 7800\, \text{kg}\, \text{m}^{-3}$, $C_{p} = 460\, \text{J}\, \text{kg}^{-1}\, {}^{\circ} \text{C}^{-1}$, and $k = 60\, \text{W}\, \text{m}^{-1}\, {}^{\circ} \text{C}^{-1}$. Determine the time required for the temperature of the bar to reach $250\, {}^{\circ} \text{C}$ by using the lumped system analysis.

\vspace{6pt}

\textbf{1--15, Solution:}

\vspace{6pt}

\textbf{1--16, Problem Statement:}
A thermocouple is to be used to measure the temperature in a gas stream. The junction may be approximated as a sphere having thermal conductivity $k = 25\, \text{W}\, \text{m}^{-1}\, {}^{\circ} \text{C}^{-1}$, $\rho = 8400\, \text{kg}\, \text{m}^{-3}$, and $C_{p} = 0.4\, \text{kJ}\, \text{kg}^{-1}\, {}^{\circ} \text{C}^{-1}$. The heat transfer coefficient between the junction and the gas stream is $h = 560\, \text{W}\,\text{m}^{-2}\,{}^\circ\text{C}^{-1}$. Calculate the diameter of the junction if the thermocouple should measure $95\%$ of the applied temperature difference in $3\text{s}$.

\vspace{6pt}

\textbf{1--16, Solution:}

\vspace{6pt}

\textbf{2--2, Problem Statement:}
A slab, $0 \leq x \leq L$, is initially at a temperature $F \left( x \right)$. For times $t > 0$ the boundary surface at $x = 0$ is kept insulated and that at $x = L$ dissipates heat by convection into a medium at zero temperature with a heat transfer coefficient $h$. Obtain an expression for the temperature distribution $T\left( x, t \right)$ in the slab for times $t > 0$ and for the heat flux at the boundary surface $x = L$.

\vspace{6pt}

\textbf{2--2, Solution:}

\vspace{6pt}

\textbf{2--6, Problem Statement:}
In a one-dimensional infinite medium, $-\infty < x < \infty$, initially, the region $a < x < b$ is at a constant temperature $T_0$, and everywhere outside this region is at zero temperature. Obtain an expression for the temperature distribution $T \left( x, t \right)$ in the medium for times $t > 0$.

\vspace{6pt}

\textbf{2--6, Solution:}

\vspace{6pt}

\textbf{2--15, Problem Statement:}
Obtain an expression for the steady-state temperature distribution $T \left( x, y \right)$ in a rectangular region $0 \leq x \leq a$, $0 \leq y \leq b$ for the following boundary conditions: the boundary at $x = 0$ is kept insulated, the boundary at $y = 0$ is kept at temperature $f \left( x \right)$ and the boundaries at $x = a$ and $y = b$ dissipate heat by convection into an environment at zero temperature. Assume the heat transfer coefficient to be the same for both boundaries.

\vspace{6pt}

\textbf{2--15, Solution:}

\vspace{6pt}

\end{document}
